\clearpage\thispagestyle{empty}
\pdfbookmark[0]{Quellen zu Auszüge aus Gauß’ Kometenbriefen vom September 1844}{cometsep1844-sources}
\begin{center}\textbf{Quellen zu Auszüge aus Gauß’ Kometenbriefen vom September 1844}\end{center}
\noindent Diese KI-Übertragung ins heutige Deutsch folgt den Auszügen vom 1., 3., 6. und 8. September 1844 auf S.~468--469 der \href{https://gdz.sub.uni-goettingen.de/id/PPN236007467}{\textit{Werke}, Band~VI} (Göttingen, 1874). Carl Friedrich Gauß schrieb an Heinrich Christian Schumacher. Die Zeitschrift druckte die vier Auszüge zusammen in den \textit{Astronomischen Nachrichten}, Band~XXII, Nr.~516, \href{https://api.digitale-sammlungen.de/iiif/image/v2/bsb10049447_00103/full/1800,/0/default.jpg}{Sp.~189--190} und \href{https://api.digitale-sammlungen.de/iiif/image/v2/bsb10049447_00104/full/1800,/0/default.jpg}{Sp.~191--192}. Die Quellenangabe der Werkausgabe datiert das Heft auf den 3. Oktober 1844.
\par\smallskip\noindent Bei der Rektaszension vom 31. August steht in der Werkausgabe ein hochgestelltes \(u\) nach der 8, im Erstdruck ein doppeltes Strichzeichen. Im Auszug vom 6. September steht in der Werkausgabe, dass das Ergebnis unter Berücksichtigung der Rektaszension beobachtet wurde. Im Erstdruck wurde es unter Berücksichtigung der Refraktion berechnet. Die Übertragung folgt den Lesarten der Werkausgabe.
\par\smallskip\noindent Gauß’ Bemerkung zum Radiusvektor bezieht sich auf die Koordinatenformeln in \href{https://api.digitale-sammlungen.de/iiif/image/v2/bsb10049447_00066/full/1800,/0/default.jpg}{Nr.~512, Sp.~116} der Zeitschrift. Dieselbe Nummer nennt die \hyperlink{native-472-notice}{Mitteilung über Goldschmidts Elemente} in der Werkausgabe. KI-Übertragung und KI-Satz.
\clearpage
